\documentclass[11pt]{article}

\usepackage[margin=1in]{geometry}
\usepackage{amsmath, amssymb, amsthm}
\usepackage{mathtools}
\usepackage[colorlinks=true, linkcolor=blue, urlcolor=blue, citecolor=blue]{hyperref}

\title{\textbf{Why gamakas resist notation}}
\author{Sanjay Suresh}
\date{January 2026}

\begin{document}
\maketitle

Western staff notation pins a note to a discrete pitch and a discrete duration. Carnatic music does
not live there --- a single swara is often a \emph{gamaka}, an oscillation or slide whose shape,
speed, and amplitude are part of the note's identity, not ornaments added on top.

So the same written phrase in two ragas can sound completely different: the notation captures the
skeleton (which swaras, in what order) but not the gesture between them, which is where the raga
actually lives. On the violin this is the whole game --- the left hand is continuously bending, and
the ``correct'' version of a note is a trajectory, not a point.

It is a good reminder that a notation is a lossy compression of the thing, chosen for what its
culture decided was worth writing down.

\end{document}
